\section{NUMBER OF PLAYERS}

    Beware! The rules of the 2–4 \faPeopleGroup\ game are very different from the ones for 5–10 \faPeopleGroup! Some
    cards have a text saying “2–4 \faPeopleGroup”: this is the side to use for games from 2 to 4 players.



\section{SET-UP}

    \begin{enumerate}
        \item Take a boss card and place it at the centre of the table. It is automated, it is
        therefore not a player.
        \item Give each player a Prolet worker card (red flag symbol, without any secret
        role) face up. They place it in front of them on the table.
        \item Give a personal token to each player. Try to find the best suited token for
        each person, so that the association will be easier to remember. Keep aside
        the hourglass token.
        \item Place the robot cards deck at the centre of the table.
        \item Remove the “Depression” Contracts card from the game. Or keep it in case
        you want to play a harder game. Place the deck of Contracts cards face
        down at the centre of the table.
        \item Place the deck of Event cards face down at the centre of the table. Assemble
        the Event cards deck with:
        \begin{itemize}
            \item 2 players: 8 random Event cards
            \item 3 players: 10 random Event cards
            \item 4 players: all 17 Event cards
        \end{itemize}
        \item Place the 6 phase cards (on the “2–4 \faPeopleGroup” side) face up, one next to the other,
        at the centre of the table. Place the phase marker token (the one with an
        hourglass drawing) on the first phase card.
        \item Create a space for the Bank where you put a pile of coins. Take plenty of
        coins, the exact amount is not important (for a 3 \faPeopleGroup\ game, you will need approximately
        30 coins, representing 30 \faRubleSign). Use some real coins, if you have
        them, for representing the \faRubleSign, the in-game currency. If you don't have coins
        at hand, you can use anything you “borrowed” from work or any other small
        and homogeneous object you can find. Alternatively, note the amounts on
        a piece of paper.
        \item All the workers and the automated boss start with 0 \faRubleSign\ coins each.
    \end{enumerate}



\section{VICTORY CONDITIONS}

    Whoever achieves the victory conditions, at any moment, can declare victory,
    and end the game. It can happen for more than one player at the same time. In
    this case, all these players win the game.

    \subsection{COLLECTIVE VICTORY}
    
        All players still in the game, win when they control as many robots as there are
        players, eliminated or not.

        After an epic victory, the players are encouraged to collect all the metallic coins
        used for playing and use them for making a donation to a workers' union or social
        movement. In the real world! We mean, out of this game, give those pennies
        some real use!
    
    \subsection{COLLECTIVE DEFEAT}
    
        All the players lose the game when:
        \begin{itemize}
            \item The automated boss has as many robots as all the players, eliminated
            and still in play;
            \item A player should draw an Event card, but the Events deck ran out of cards;
            \item Half of the players are eliminated (when playing in 3, this happens with the
            elimination of 2 \faPeopleGroup).
        \end{itemize}
    
    \subsection{CLASS BETRAYAL VICTORY}
    
        Any non-Stressed player with 8 \faRubleSign\ coins, or two non-Stressed players united by
        the Marriage Event card with 10 \faRubleSign, can declare themselves the winner(s) as class
        traitor(s)!
        If more than one player fulfils the Class Betrayal victory conditions at the same
        time, they can all declare victory.



\section{GAMEPLAY}

    The game is played in rounds until 1 or more players achieve the Class Betrayal
    victory conditions or the workers' collective victory or collective defeat conditions
    are met. When this happens, the game immediately ends.
    
    A round consists of 6 phases. The starting player changes every round and is in
    charge of revealing one Event and one Contracts card. Each phase is played differently.
    
    The first starting player is the one who most recently participated in a protest
    or strike (in real life or in the game).

    \begin{itemize}
        \item Phase 1 State-funded innovation: the boss could buy a robot.
        \item Phase 2 Event: the starting player reveals an Event card and applies its effect.
        \item Phase 3 Contracts: the starting player reveals a Contracts card, and picks
        a Contract first. Clockwise, the other players choose between the remaining
        Contracts.
        \item Phase 4 Strike: simultaneously, non-Stressed workers reveal whether they
        will strike or work. Players cannot discuss strike intentions.
        \item Phase 5 Payday: players who worked get a salary. Robots produce.
        \item Phase 6 Rent: the players pay the Rent to the automated boss. The player
        on the left of the starting player is the new starting player.
    \end{itemize}



\section{Phase 1 STATE-FUNDED INNOVATION}

    \subsection{AUTOMATED BOSS ACTIONS}
    
        \begin{itemize}
            \item If the automated boss has at least 15 \faRubleSign, move 15 \faRubleSign\ to the Bank, and place
            next to the boss card a new robot card on the blue side.
        \end{itemize}



\section{Phase 2 EVENT}

    \subsection{WORKERS ACTIONS}
    
        \begin{itemize}
            \item The starting player must reveal an Event card and apply its effect. If there is
            no card left to pick, see collective defeat.
        \end{itemize}

    \subsection{ADDITIONAL DETAILS AND EXAMPLES}

        \subsubsection{SYMBOLS ON THE EVENT CARDS}

            PLAY AND DISCARD

            PERMANENT EFFECT, the affected player keeps the card

            This card affects the CURRENT STARTING PLAYER

            This card affects a CHOSEN PLAYER

            This card affects MORE THAN ONE PLAYER

            If the unavoidable and immediate effect of an Event card is to eliminate a player,
            ignore this effect: Event cards cannot directly eliminate players.

        \subsubsection{EXAMPLE}

            \textit{The automated boss has no coins at the beginning of the game, so clearly will not buy a
            robot just yet. Between the 3 workers, Asdrubale, Bedelia and Candida, the last participating
            in a protest was Bedelia (the management of the data centre where she works is blocking
            the negotiation of the collective agreement, so she plugged a bunch of short network cables
            to adjacent ethernet ports of many switches), so she starts the game revealing the first Event
            card: the “Gambling” card says “one of the players with most coins gives all but 1 coin to
            the boss”, but none of the 3 \faPeopleGroup\ have any coins, and the card is discarded (and no one gets
            Stressed).}



\section{Phase 3 CONTRACTS}

    \subsection{WORKERS ACTIONS}

        \begin{itemize}
            \item The starting player must reveal a Contracts card, choose a Contract and
            place the personal token on it. The other players, in clockwise order, must
            choose between the remaining Contracts listed on the card, placing their
            personal token on the chosen Contract.
        \end{itemize}
    
    \subsection{ADDITIONAL DETAILS AND EXAMPLES}

        \begin{itemize}
            \item Get Stressed: Contracts may cause Stress, this will be applied during phase
            5 Payday.
            \item Holidays included: in addition to the salary, Contracts may include Holidays.
            Holidays will allow players to eliminate Stress during phase 5 Payday.
            \item Permanent Contract: will last for more rounds, until the player accepts another
            Contract or participates in a strike. This player will keep this card next
            to the character card during the Contract's validity.
            \item A Contracts card cannot directly eliminate a player: if a Contract says “get
            Stressed” but you are already Stressed, you are unaffected.
            \item When the Contracts cards run out, shuffle the discarded Contracts cards (except
            the permanent Contract cards still in use by the players).
        \end{itemize}

        \subsubsection{SYMBOLS ON THE CONTRACTS CARDS}

            Get Stressed: Contracts may cause Stress, this will be applied during
            phase 5.

            Holidays included: in addition to the salary, Contracts may include Holidays.

            Holidays will allow players to eliminate Stress during phase 5.

        \subsubsection{EXAMPLE}

            \textit{Since it's her turn to start, Bedelia draws a Contracts card, and shows it to everyone. It's the
            “Boom” card! She chooses the “3-coins Contract (with Stress)” first. On her left, Candida,
            picks the “2-coins Contract” from the remaining Contracts. Finally, Asdrubale takes the “1-coin
            permanent Contract (until strike)”. Asdrubale then places the Contracts card next to his
            Prolet character card, keeping it until he decides to rescind the Contract. On the card, there
            is a fourth Contract that no one will take, as there are only 3 players.}



\section{Phase 4 STRIKE}

    \subsection{WORKERS ACTIONS}

        \begin{itemize}
            \item All the non-eliminated non-Stressed players place one hand towards the
            centre of the table and, at the same time, must vote one of the following:
            \begin{itemize}
                \item thumb up = strike for Acceleration
                \item closed fist = strike for Expropriation
                \item open hand = work
            \end{itemize}
        \end{itemize}
Alternatively, they can write their action on a piece of paper.

    
    \subsection{ADDITIONAL DETAILS AND EXAMPLES}

        \begin{itemize}
            \item Stressed players cannot strike and so must work.
            \item It's forbidden to talk or communicate about strikes.
            \item Acceleration Strike: every thumb raised creates one robot for the boss.
            The boss loses all the coins earned up to that point. Beware that more than
            one robot can appear in the same round! For example, if all the workers
            vote raising their thumbs, the boss loses all the savings and immediately
            receives those many robots, triggering the collective defeat!
            \item Expropriation Strike: If the majority of players (with 2 \faPeopleGroup\ the majority is 2;
            with 3 \faPeopleGroup\ majority is 2; with 4 \faPeopleGroup\ majority is 3) voted with a closed fist, take
            one robot from the boss, flip it to the red side and assign it to a single player.
            You can also expropriate a robot created by an Acceleration Strike in this very
            round. If the strike does not reach the majority, nothing happens, but the
            striking players will not get any salary in phase 5 Payday.
        \end{itemize}
    
        \subsubsection{EXAMPLE}

            \textit{After a short countdown, the 3 \faPeopleGroup\ express their decisions with their hands: Bedelia
            wants to work (secretly intends to win with a Class Betrayal) and so keeps her
            hand open; Candida also intends to work, and therefore exposes her open hand;
            Asdrubale wasn't really interested in the permanent Contract, and decided to do
            an Acceleration Strike! He is showing his thumb up, which indicates the Acceleration
            Strike. Consequently, he returns the permanent Contract card to the area of
            discarded Contracts, and places a new blue robot next to the boss. The automated
            boss discards all the owned coins, if any.}



\section{Phase 5 PAYDAY}

    \subsection{AUTOMATED BOSS ACTIONS}

        \begin{itemize}
            \item Each robot the boss controls produces 2 \faRubleSign\ for the boss.
        \end{itemize}
    
    \subsection{WORKERS ACTIONS}

        \begin{itemize}
            \item Each non-striking worker receives the salary and applies the additional
            effects from the chosen Contract. The effect of “with Holidays” is to remove
            the Stress, Untapping the character card. The effect of “get Stressed” is to
            Tap the character card (remember that a Stressed player cannot do any kind
            of strike).
            \item Each Expropriated robot produces 1 coin for the worker it is assigned to.
        \end{itemize}

    \subsection{ADDITIONAL DETAILS AND EXAMPLES}

        \begin{itemize}
            \item Players who attempt any kind of strike will receive neither the salary from
            the chosen Contract, nor its positive or negative effect (i.e. “get Stressed” or
            “with Holidays”).
        \end{itemize}
    
        \subsubsection{EXAMPLE}

            \textit{The new robot produces 2 \faRubleSign\ for the boss. Bedelia receives 3 \faRubleSign\ from the Bank, but
            has to Tap her Prolet character card, as the chosen Contract indicated that she
            would have become Stressed. Candida receives 2 \faRubleSign\ from the Bank. Asdrubale went
            on strike, so he receives nothing.}



\section{Phase 6 RENT}

    \subsection{WORKERS ACTIONS}

        \begin{itemize}
            \item Unless an Event card changes this, each worker pays 1 coin directly to the
            boss.
        \end{itemize}
    
    \subsection{ADDITIONAL DETAILS AND EXAMPLES}

        \begin{itemize}
            \item A player who can't pay the Rent gets Stressed. If already Stressed and
            cannot pay the Rent, this player is considered eliminated and stops participating in the game.
            \item If a player needs to pay Rent of 2 \faRubleSign\ (e.g. due to the “Birth” Event card), but
            has only 1 coin, they pay nothing but become Stressed. If they are already
            Stressed, they are eliminated.
            \item If half of the players are eliminated, the game is over, see collective de-
            feat. In a game with 3 \faPeopleGroup, this happens when 2 \faPeopleGroup\ are eliminated.
        \end{itemize}
        The player on the left of the starting player is the new starting player. The game
        continues with phase 1, and follows all the phases again.


        \subsubsection{EXAMPLE}

            \textit{Bedelia pays her 1-coin Rent, moving one of her 3 \faRubleSign\ next to the boss card. Candida
            does the same thing. Asdrubale, having no coins for paying the Rent, has to Tap his
            Prolet card (he is now Stressed).}



\section{ADDITIONAL RULES}

    \begin{itemize}
        \item An elimination can happen only if a player cannot pay the Rent during phase
        6 and is already Stressed.
        \item If some effect of an Event card cannot be applied, skip it and apply the remaining
        effects on that card.
        \item It is forbidden to communicate about strikes: only the voting during phase 4
        will reveal the workers' actions.
        \item Workers cannot keep their coins hidden.
        \item Workers pay the Rent directly to the automated boss.
        \item In case of disagreement between players (e.g. who should control an expropriated
        robot), the starting player will decide.
    \end{itemize}



\section{NOTES ON THE EVENT CARDS}

    \paragraph{Birth}
        If a married couple (see Marriage card) has a Birth, they can decide who
        takes care of the child in case there will be a Divorce Event.
    \paragraph{Collective Agreement}
        Each Contract on this round's Contracts card can be
        chosen by more than one player.
    \paragraph{Dentist}
        If possible, one player should pay or get Stressed. If no player has 2 \faRubleSign,
        and all players are already Stressed, this card is the least of your problems and
        you can ignore it.
    \paragraph{Divorce}
        If the Marriage card is in game, both of the involved players lose all the
        coins, and remove the Stress, if not, ignore this card.
    \paragraph{Eviction}
        The card effect is applied to the player who has the Gentrification card.
    \paragraph{Gentrification}
        Once one player has this card, the only way to get rid of it is that
any player draws the Eviction card.
    \paragraph{Mandatory Overtime}
        If the player who picked this card participates in a strike
        this round, they do not get the 1 coin extra, and do not get Stressed due to this
        card.
    \paragraph{Marriage}
        You can bring the characters cards closer together and put this card
        in the middle. Even if married, the Rent is 1 coin for each worker. What changes
        is that one player may pay both Rents. Marriage modifies the Class Betrayal winning
        conditions, check them out.
    \paragraph{Production Bonus}
        If the player who picked this card participates in a strike this
        round, they do not receive the extra 2 \faRubleSign.
    \paragraph{Rotating Shifts}
        Participants in a strike this round, do not get the 1 coin extra,
        and do not get Stressed due to this card.
    \paragraph{Sabotage}
        If there is only a red robot (an expropriated one), destroy it, if there
        are no robots in play, ignore this card.
    \paragraph{Workers' Union}
        If the boss is out of money, you are out of luck.



